\documentclass[a4paper,11pt]{article}
% !TeX program = pdflatex
\pdfoutput=1
\pdfminorversion=7
\usepackage[utf8]{inputenc}
\usepackage{jheppub}

\usepackage{amsmath,amssymb,latexsym}
\usepackage{braket}
\usepackage{bm}
\usepackage{booktabs}
\usepackage{tabularx}
\usepackage{longtable}
\usepackage{cancel}
\usepackage[svgnames]{xcolor}
\usepackage{hepunits}
\usepackage{mathtools}
\usepackage{graphicx}
\usepackage[most]{tcolorbox}
\usepackage{caption}
\usepackage{tikz}
\usetikzlibrary{calc}
\def\ie{\begin{equation}\begin{aligned}}
\def\fe{\end{aligned}\end{equation}}
\newtcolorbox{proofbox}{
  enhanced, breakable, colback=gray!8, colframe=black,
  boxrule=0.8pt, arc=0pt, outer arc=0pt,
  left=8pt, right=8pt, top=8pt, bottom=8pt,
  before skip=10pt, after skip=10pt
}
\hypersetup{colorlinks=true,linkcolor=blue,citecolor=blue,urlcolor=blue}
\setlength{\emergencystretch}{3em}
\makeatletter
\def\verbatim@font{\normalfont\ttfamily\small}
\makeatother

\title{dSIBP Handbook}

\begin{document}
\sloppy
\setcounter{tocdepth}{1}
\maketitle

\section{Scope and quick start}

dSIBP is a topology-driven Wolfram Language package for de Sitter Feynman integrals.
It constructs time and loop-momentum IBP seeds, parameter-derivative operators, sector
metadata, backend-neutral \texttt{linearData}, and Kira input files. It does not perform
reduction internally. External reducers such as Kira or Rational Tracer remain separate
steps controlled by the user.

The package supports arbitrary loop numbers and topologies, mixed massive and massless
propagators, exact contact-sector construction, EOM and symmetry canonicalization,
user-defined master bases, and differential-equation and scaling checks. The public
topology schema is deliberately strict: package metadata is derived from the input
graph rather than inferred from symbol names.

A standard session loads the package, initializes a family, generates symbolic seeds,
expands continuous indices, and constructs linear data:
\begin{verbatim}
packageDir = FileNameJoin[{Directory[],"dSIBP"}];
If[!MemberQ[$Path,packageDir],AppendTo[$Path,packageDir]];
Needs["dSIBP`"];

context = DSInit[case];
seedData = DSSeeds[context];
allSeeds = DSAllSeeds[seedData];
ibpData = DSGenerateIBP[allSeeds,{-1,1}];
linearData = DSLinear[ibpData,context];
\end{verbatim}
The package can be installed at any path; loading requires only the standard module tree and does not depend on the directory depth.
Use \texttt{DSKiraPlan} and \texttt{DSKiraExport} when Kira input is required, run Kira
outside the package, and use \texttt{DSKiraImport}, \texttt{DSDE}, and
\texttt{DSScaleCheck} only after complete external results are available.

The bundled examples cover six complementary workflows:
\begin{itemize}
  \item \texttt{01\_mixed\_bubble\_workflow.wl}: basic initialization, generation,
  linearization, integral ordering, and planning;
  \item \texttt{02\_function\_system\_hankel.wl}: Hankel and related function-system input;
  \item \texttt{03\_single\_massive\_sunrise}: general two-loop seeds and parameter operators,
  without sampled relations or reduction;
  \item \texttt{04\_pure\_massive\_bubble\_closed\_loop}: a 19-master Kira/DE/reference/scaling
  closed loop;
  \item \texttt{05\_tree\_two\_vertex\_time\_ibp}: direct pure-time/tree relations and
  naive versus dlog differential equations;
  \item \texttt{06\_mix\_bubble\_tree}: mixed cycle/bridge propagators, custom coordinates,
  contact sectors, and an 81-master reference summary.
\end{itemize}

\section{A complete mixed bubble-tree example}

A useful nontrivial input contains a massive cycle propagator and a massless cycle
propagator between $v_1$ and $v_2$, followed by a massless fixed bridge from $v_2$ to $v_3$.
Only the two cycle propagators carry loop-integral parity constraints. The bridge has no
continuous loop-power slot and is excluded from parity.

Momentum declarations serve two distinct roles. \texttt{loopExternalMomenta} specifies
independent directions entering the loop, while \texttt{independentExternalMomenta}
specifies the magnitudes of momenta that remain outside it. The package constructs the
required Gram coordinates and checks whether these declarations are complete, overcomplete,
or undercomplete. It never infers a role from a symbol's name.

Initialization returns the parsed topology, contact-reachable sectors, canonical packs,
zero points, normalization data, and capability reports. A typical high-level sequence is
\begin{verbatim}
context = DSInit[caseInput];
seedData = DSSeeds[context];
allSeeds = DSAllSeeds[seedData];
generated = DSGenerateIBP[allSeeds,
  {a[v1],-1,3},{a[v2],-1,3},{b[1],-2,3},{b[2],-2,3}];
linearData = DSLinear[generated,context];
\end{verbatim}
The same example also demonstrates custom coordinates through
\texttt{DSRedefineParameters}, fixed-bridge magnitudes, sector contact maps, and exact
round trips between user and internal scalar products.

\section{From SK rules to integral families}

\subsection{Vertices and external phases}

Each vertex is an association containing \texttt{id}, \texttt{vertexType}, and
\texttt{externalLegEnergy}. The vertex type fixes the sign of the external phase, and the
energy is the independent phase variable. The parser normalizes these signs once; all
subsequent time derivatives and contact rules use the same convention.

\subsection{Propagator types}

The propagator classes are massive full, massless full, massive cross, and massless
cross or fixed lines. Massive full lines have a two-theta contact system and a Hankel-type
function space, while massless full lines use a two-vertex quotient representation. Cross
lines have no ordinary two-theta shrink channel. Fixed tree lines carry external magnitudes
but no continuous loop-integral power.

A standard line uses \texttt{id}, \texttt{endpoints}, \texttt{momentum}, and
\texttt{massType}. Massive lines additionally require \texttt{nu}. A custom massive
function system may be supplied through \texttt{functionSystem}; no physical type is
inferred from a mass parameter alone.

\subsection{Contact-reachable sectors}

The root topology determines loop number, cycle/bridge roles, and propagator schemas.
Contact operations only shrink lines, merge vertex representatives, and update zero points
and normalization. They do not create a new unrelated family or silently reduce the
structural loop number.

Massive contact terms are compiled from the Wronskian system. Simultaneous contacts merge
a vertex only once, while integer shifts from all selected lines are accumulated. Uncut
propagators are canonicalized in the target sector. Every reachable sector retains the
metadata required by seeds, derivatives, and serializers.

\section{Integral representations}

In full mode the public integral head is
\begin{verbatim}
J[aList,linePacks,ispList]
\end{verbatim}
The first slot contains compact vertex-time powers. Each full-line pack has the form
$\{b_e,n_{e,1},n_{e,2}\}$, where $b_e$ is a continuous denominator power and the two
entries are discrete building-block states. A shrunk line uses the corresponding compact
$bS$ representation. ISP powers are collected in the third slot and are measured from zero.
Positive ISP powers are numerators; explicitly requested negative powers are also allowed.

In \texttt{ibpMode->"timeOnly"}, the public representation is instead
\begin{verbatim}
J[sectorKey,timeShifts,stateBits]
\end{verbatim}
Here \texttt{sectorKey} is a fixed-width string of full/shrunk bits in root-line order.
The other two slots contain compact time shifts and the sector-local discrete-state registry.
The two schemas are intentionally distinct; no adapter silently rewrites full-loop Kira
artifacts into the time-only representation.

\section{Massive function systems, derivatives, and EOM}

A massive building block is described by a $2\times2$ first-order system. Its derivative
matrix and Wronskian determine, respectively, derivative shifts and contact-shrink shifts.
These two data channels are not interchangeable.

The \texttt{DSInit} option \texttt{NuConvention} selects the writing of every baseline
formula that involves \texttt{nu}. The default \texttt{NuConvention->"Positive"} uses the
input Hankel order as supplied; \texttt{NuConvention->"Negative"} flips its sign per massive
line and byte-reproduces the baseline equations of earlier releases. The setting is frozen
per line at initialization, and no later command accepts a convention option.

After a generator acts, any Hankel state of second order or higher is immediately reduced
by the equation of motion. This happens before linear-system construction, so no symbolic
$n\geq2$ state is passed to a backend. Contact shifts use the compiled Laurent coefficients
and preserve the associated zero-point and prefactor changes.

Parameter derivatives act on both the integral and its explicit coefficients. The package
also accounts for sector normalization, so a derivative of $J_s=N_s I_s$ includes the
$D(\log N_s)$ term and, across a contact, the appropriate $N_s/N_t$ ratio.

\section{Massless quotients and time IBP}

Massless full propagators retain an ordered two-vertex three-slot pack. Regular and delta
terms are handled in that same representation, including contact shrink and coincident
cancellation. Several full lines joining one vertex pair are treated by a common-theta
construction rather than by independently duplicated contact rules.

An odd subset of lines in a shared bundle determines the contact transition. Its coefficient
and target-sector mapping are evaluated once, and duplicate transitions are removed by source,
target, and selected-line identity. Massless fixed lines do not acquire a fictitious loop
parity constraint.

Parallel massless full lines joining the same vertex pair are additionally folded by a
built-in canonical layer. In the original per-line index representation the product kernel of
such a bundle depends only on the total parity $\sum_en_e\bmod 2$ of the endpoint states, so
every member is mapped to the representative in which the first line in line order carries
that parity and all remaining lines are zero; a member whose endpoint order is reversed
contributes $(-1)^{n_e}$ per line. The fold sits after the single-line quotient and before
symmetry in the canonical chain.

Seed enumeration distinguishes two cases. Without explicit parity constraints, or for a
bundle no constraint touches, every member of the bundle is still enumerated over $00/10$:
one-sided states carry subsector boundary information that the fold only normalizes later.
When a user constraint in \texttt{"parityConstraints"} touches a member of a bundle (banana
total-parity sampling), initialization sums the constraints touching that bundle (including
$b$ indices, mod $2$) into its total-parity constraint
\begin{equation}
 \sum_{e\in\mathcal B}\bigl(b_e+n_{e,1}+n_{e,2}\bigr)\equiv r_{\mathcal B}\pmod 2,
 \label{bundleTotalParityManualEn}
\end{equation}
and sampling enumerates only the parity slot $\{0,1\}$ of the representative line (smallest
line index), with the remaining members fixed to $0$. \texttt{DSInit} prints a bilingual red
notice listing which bundles are merged and the resulting total-parity target of each, then
continues: it neither rejects nor stops. If only some members are constrained, the same rule
sums the provided constraints and the resulting subsystem may leave the originally selected
one; several parity subsystems may be selected simultaneously, and propagators without a
parity selection are sampled over all parities.

A massless full line whose two endpoints coincide is converted to the contracted single-slot
pack $\{b_{S,e}\}$, sharing all bookkeeping with shrunk lines. This covers lines declared with
equal \texttt{endpoints} and lines that become coincident when a contact merges their vertices.
At $\Delta=0$ the kernel blocks take $M_0=2\theta(0)=1$ and $M_1=0$, so the undifferentiated
block is the constant $1$, antisymmetric states vanish identically, and only the residual
momentum-magnitude power survives. Declared self-loops convert at initialization with
$b_{S0,e}=b_{0e}$ and no endpoint-state labels; the root graph audit still counts them as
active cycle lines, since the residual $b$ power carries the line's own loop momentum. In
sector construction, coincident survivors join the contracted set of the sector, and distinct
contact routes that reach the same contracted set define a single sector entry, while contact
routing still follows each individual selection. Odd-state integrals are zeroed in the massless
endpoint canonicalization against the source topology, and even-state contact packs
$\{b,0,0\}$ convert to the single slot automatically. A massive self-loop is not converted: it
keeps its three-slot pack and the built-in reversal rule, and converted single-slot self-loops
still enter the odd-ISP zeroing of that rule.

The direct pure-time/tree path uses the same physical contact data but presents them in a
sector-tagged public representation. Formula-based tree reduction is available only for
massive-only quotients whose derivation has been certified. In all other cases, the package
returns a structured \texttt{PendingRederivation} result.

\section{Topology, scalar products, ISP, and symmetry}

The required topology input includes \texttt{vertices}, \texttt{lines},
\texttt{loopMomenta}, \texttt{loopExternalMomenta}, and
\texttt{independentExternalMomenta}. ISP definitions are required when the propagators do
not span all independent scalar products. Every ISP must be closed under the selected loop
momentum generators.

The package distinguishes loop-external directions from loop-free magnitudes. Default loop
coordinates are $\texttt{ssij}=\sqrt{\texttt{sp}[k_i,k_j]}$, with the derivative convention
\begin{equation}
 \frac{\partial}{\partial\texttt{ssij}}
 =2\,\texttt{ssij}\frac{\partial}{\partial\texttt{sp}[k_i,k_j]}.
 \label{handbookRootDerivative}
\end{equation}
Custom coordinates are introduced through \texttt{KinematicRules} or
\texttt{DSRedefineParameters}. Undercomplete declarations stop initialization;
overcomplete declarations may permit symbolic seed generation but disable unique
derivatives and backend export.

Symmetry rules are applied as directed replacements. Users must choose one canonical
representative for each equivalence class and provide oriented rules only. The symmetry of
\texttt{sp} itself does not imply a graph symmetry: masses, energies, and momentum routing
must also match.

\section{IBP generators and atomic operations}

Time generators differentiate with respect to conformal times. Momentum generators include
loop scalings and the independent external directions required by the declared loop basis.
The complete set is constructed before any index sampling.

The public atomic operations include \texttt{dtau}, \texttt{dqq}, \texttt{dqk}, and
\texttt{ds}. The first three generate time and momentum IBP relations, while \texttt{ds}
is the total parameter derivative used by the differential-equation workflow. Conversion
helpers \texttt{rep2innerform} and \texttt{rep2outform} translate between user and internal
coordinates without changing integral indices. \texttt{rep2Integrand} exposes an inert
integrand for inspection; it does not integrate or generate a relation.

\section{Seeds and continuous-index envelopes}

\subsection{Symbolic templates and shift metadata}

\texttt{DSSeeds} enumerates the independent discrete templates before continuous sampling.
Massive vertices use the four $0/1$ states; massless quotients use their proved canonical
representatives, with parallel bundles folded to the first-line parity form. EOM, symmetry,
and sector canonicalization are applied while continuous
indices remain symbolic.

\texttt{DSAllSeeds} returns the flat template list, while \texttt{DSSeedGroups} and
\texttt{DSSeedGroupMetadata} expose the recommended grouping. \texttt{DSMetaSeedRange}
records the affine shift set of every continuous index. A flat input is treated as one
group; each top-level element of a nested input is treated as a separate group.

A uniform target envelope is written as
\begin{verbatim}
ibp = DSGenerateIBP[allSeeds,{L,U}];
\end{verbatim}
A detailed envelope assigns an interval to every root index. For an index with shift set
$\Delta$, the corresponding seed interval is the intersection of all inverse images:
\begin{equation}
 x\in[L-\min\Delta,\;U-\max\Delta].
 \label{handbookSeedInverse}
\end{equation}
For ISP indices, the lower endpoint is also bounded below by the user target lower endpoint.
Empty inverse ranges are reported explicitly and disable the relevant completeness claim.

\subsection{Seed-envelope cutting}

\texttt{SeedEnvelopeCut} filters the inferred seed points before parity filtering and
template substitution. The following calls disable cutting, use the default, and provide
custom parameters:
\begin{verbatim}
ibp0 = DSGenerateIBP[allSeeds,{-1,1},SeedEnvelopeCut->None];
ibp1 = DSGenerateIBP[allSeeds,{-1,1},SeedEnvelopeCut->Automatic];
ibp2 = DSGenerateIBP[allSeeds,{-1,1},SeedEnvelopeCut->{10,4,5,4}];
\end{verbatim}
Let $C_+$ be the sum of positive continuous slots, $C_-$ the absolute sum of negative
slots, and $T=C_++C_-$. For $C_+\geq C_-$, acceptance requires $T\leq Ex_1$ and
$C_-\leq Ex_2$; otherwise it requires $T\leq Ex_3$ and $C_+\leq Ex_4$. The default for a
group of $N$ continuous slots is $\{N+3,4,4,4\}$. Every item of an explicit four-tuple may
be an integer or \texttt{Automatic}, in which case only that default item is substituted.

The returned \texttt{seedEnvelopeCut} metadata records the resolved parameters, accepted
and skipped point counts, and groups emptied by the cut. \texttt{None} restores the uncut
point domain.

\subsection{Canonical-support cutting}

\texttt{TotalPowerLimit} is an independent filter on canonical equation support:
\begin{verbatim}
ibp3 = DSGenerateIBP[allSeeds,{-1,1},TotalPowerLimit->7];
ibp4 = DSGenerateIBP[allSeeds,{-1,1},
  SeedEnvelopeCut->{10,4,5,4},TotalPowerLimit->7];
\end{verbatim}
For an integer $m$, the package first combines equal integrals and removes zero coefficients.
It then retains an equation only if every canonical integral in its actual support lies in
the per-slot target box and has continuous-index absolute sum at most $m$. The default
\texttt{None} disables this filter. The two cut options can be used independently or together.

The full-mode fast path decomposes Association-based sealed templates into integral
skeletons, integer shifts, and coefficient pieces, then memoizes canonical results by raw
integral. Raw expression lists and failed templates use the original symbolic path. This
changes performance, not the meaning of either cut option.

These options do not prove closure for an arbitrary target or basis.
\texttt{completeSystemQ} describes generation completeness within the selected envelope
and cutting policy. A new target, basis, or cut setting still requires an external
reduction check that every requested target reduces to the requested masters.

\section{Linear data and the Kira workflow}

\texttt{DSLinear} converts generated relations into backend-neutral \texttt{linearData}.
Its \texttt{integralList} is the authoritative backend order, and equations are deduplicated
by their mathematical coefficients, constant terms, and nonlinear terms rather than by
their provenance. \texttt{DSReorderIntegrals} explicitly changes that order before basis
or backend construction. \texttt{DSUserMI} attaches a reversible ordered basis of single
integrals or homogeneous linear combinations and constructs analytic derivative targets.

A formal Kira workflow has two stages: optional pre-reduction for basis exploration and a
formal stage using the active basis and its minimal derivative closure. \texttt{DSKiraPlan}
constructs the plan, and \texttt{DSKiraExport} serializes equations, maps, jobs, and manifest
information without running Kira.

For a massive bubble, \texttt{numericStage->"postDerivative"} substitutes the declared
parameter rules only after the analytic derivatives have been constructed; it does not
necessarily make the reduction purely numerical. In the bundled example, \texttt{dim},
\texttt{nu}, \texttt{etaNu}, and \texttt{analyticRegulator} take exact rational values,
but the kinematic variables remain symbolic during reconstruction. Kira therefore
reconstructs in \texttt{dsc2} and \texttt{ip0}, and the resulting DE matrices remain
symbolic in \texttt{ss11} and $P_0$. Comparison with the analytic reference at a fixed
numerical point is a separate final check.

After external reduction, \texttt{DSKiraImport} checks manifest identity, target coverage,
master order, and RHS closure. \texttt{DSDE} constructs matrices and sources in the active
basis, and \texttt{DSScaleCheck} checks the applicable symbolic, fixed-section, or exact-point
scaling certificate. A global residual master outside the requested target closure is not
by itself a failure; acceptance is based on the requested targets and their RHS.

\section{Initialization, messages, and inspection}

\texttt{DSInit} validates topology, ISP closure, kinematics, and sector metadata.
\texttt{DSInfo} returns initialization summaries. \texttt{DSKinematics} and
\texttt{DSParameterNotation} expose coordinate proposals and current notation, while
\texttt{DSRedefineParameters} performs a complete reinitialization under new rules.

\texttt{DSMessagesOn}, \texttt{DSMessagesOff}, and \texttt{DSMessagesQ} control optional
information and warnings. Fatal errors remain visible. \texttt{AuditLevel} controls whether
cached producer certificates are reused or full development audits are recomputed.
\texttt{DSPublicAPI} returns the public function groups and the high-level options actually
present in the loaded package.

\section{Pure massive bubble closed loop}

The bundled two-vertex massive bubble uses the $--$ branch, even parity on both cycle
propagators, and equal vertex energies. Its target envelope is $a_v\in[0,3]$ and
$b_e\in[-3,5]$, covering the 19-element basis and all derivative support. The example uses
\texttt{SeedEnvelopeCut->Automatic} and \texttt{TotalPowerLimit->None}.

A fresh run generates 20228 IBP relations before linearization, followed by 2442 canonical
relations, 2460 formal equations, and 1317 source integrals. Kira reduces all 300 requested
targets, leaving none unreduced. The active basis has 19 masters; five additional global
masters do not appear on the right-hand sides of the requested reductions. Import,
$19\times19$ differential equations, and scaling pass 23 checks. The physical $P_0$,
backend \texttt{ip0}, and physical $k_s$ matrices each agree with the frozen analytic
reference in all 361 entries.

The example requires an external workspace selected by \texttt{DSIBP\_KIRA\_WORKSPACE}; it must lie outside the package's Git repository (or, when the package is not inside a Git repository, outside the delivery directory), and the script refuses a workspace inside it.
The package writes inputs but does not start Kira. After Kira finishes,
\texttt{post\_kira\_check.wl} reads the existing results without regenerating inputs.
Repository files retain only lightweight input and result summaries, not databases,
reduction tables, caches, or runtime logs.

\section{Tree vertex families}

The direct tree workflow exposes \texttt{DSTreeSeeds}, \texttt{repIterative},
\texttt{DSTreeNaiveIBP}, \texttt{DSTreeNaiveDE}, and \texttt{DSTreeDLogDE}. It constructs
sector-tagged relations, ordered masters, normalized differential equations, and scaling
checks from the same contact and normalization metadata.

\texttt{DSTreeNaiveIBP} solves one-step raising relations in a specified basis.
\texttt{DSTreeNaiveDE} uses that reduction to construct the naive connection, while
\texttt{DSTreeDLogDE} constructs the block-structured dlog connection. For massless full
quotients, formula-based reduction is available only after the required identities have
been rederived.

\section{Common errors and current boundaries}

Malformed topology, missing momentum directions, incomplete ISP closure, unknown or
repeated range indices, symbolic discrete indices, and invalid cut specifications are
reported as structured failures. Empty seed ranges are not silently ignored. Backend export
also rejects unresolved internal coordinates, unsupported complex conventions, stale
manifests, and incomplete reduction artifacts.

dSIBP does not run reduction internally. It does not infer physical symmetry from symbol
names, and it does not turn a finite-envelope completeness claim into a proof for arbitrary
targets. Formula-based tree reduction remains limited to the certified massive-only
quotient cases. Reference summaries record their evidence baseline and are not silently
reused as proof for a different generation policy.

\section{Public API summary}

\small
\begin{longtable}{@{}p{0.46\textwidth}p{0.46\textwidth}@{}}
\toprule
\textbf{Function group} & \textbf{Purpose} \\
\midrule
\texttt{DSInit}, \texttt{DSInfo}, \texttt{DSKinematics},
\texttt{DSParameterNotation}, \texttt{DSRedefineParameters}
& Initialization, topology and coordinate inspection, and redefinition. \\
\texttt{DSSeeds}, \texttt{DSAllSeeds}, \texttt{DSSeedGroups},
\texttt{DSSeedGroupMetadata}, \texttt{DSMetaSeedRange}
& Symbolic discrete templates, grouping, and shift metadata. \\
\texttt{DSGenerateIBP}
& Continuous expansion with \texttt{SeedEnvelopeCut} and \texttt{TotalPowerLimit}. \\
\texttt{DSLinear}, \texttt{DSReorderIntegrals}, \texttt{DSUserMI}
& Backend-neutral equations, authoritative order, and active bases. \\
\texttt{DSKiraPlan}, \texttt{DSKiraExport}, \texttt{DSKiraImport}
& Kira planning, serialization, and verified result import. \\
\texttt{DSDE}, \texttt{DSScaleCheck}
& Differential equations and scaling certificates. \\
\texttt{DSTreeSeeds}, \texttt{repIterative}, \texttt{DSTreeNaiveIBP},
\texttt{DSTreeNaiveDE}, \texttt{DSTreeDLogDE}
& Direct pure-time/tree generation and reduction. \\
\texttt{dtau}, \texttt{dqq}, \texttt{dqk}, \texttt{ds}
& Atomic time, momentum, and parameter derivatives. \\
\texttt{rep2innerform}, \texttt{rep2outform}, \texttt{rep2Integrand},
\texttt{symmetry}, \texttt{repSymmetry0}
& Coordinate conversion, inert integrands, and symmetry access. \\
\texttt{DSMessagesOn}, \texttt{DSMessagesOff}, \texttt{DSMessagesQ},
\texttt{DSPublicAPI}
& Message control and public-interface inspection. \\
\bottomrule
\end{longtable}
\normalsize

The exact options of a function are those returned by \texttt{Options[function]}.
\texttt{DSGenerateIBP} exposes \texttt{SeedEnvelopeCut} (default \texttt{Automatic})
and \texttt{TotalPowerLimit} (default \texttt{None}), together with the shared audit
and progress options.

\end{document}
